% !TeX root = ../深度学习与大语言模型技术原理.tex
% 由主文档通过 \input 载入；请编译主文档。
\setcounter{chapter}{20} % 保留原书章号。
\chapter{机制可解释性与因果干预}
\label{ch:interpretability}

\section{从相关激活到计算机制}

机制可解释性（mechanistic interpretability）研究模型内部哪些表示和计算路径参与了某种行为。激活、神经元、注意力头和电路（circuit）是分析对象，但单个神经元可能响应多个概念，一个概念也可能分布在多个方向。注意力权重显示信息聚合方式的一部分，不能单独作为完整决策解释。

探针（probing）训练简单模型从隐藏表示预测某属性。若探针成功，说明表示中存在可读取信息；这不证明原模型实际使用该信息，也不说明该属性有因果作用。探针容量过大可能自己学习复杂规律，因此应有随机标签、不同层和容量控制等对照。

\section{激活替换与因果追踪}

激活替换（activation patching）比较干净输入和受扰输入，将某层某位置的激活从干净运行复制到受扰运行，再观察目标输出是否恢复。因果追踪（causal tracing）利用类似干预定位信息影响路径。实验必须说明扰动如何构造、替换什么张量以及采用何种输出指标。

若替换恢复行为，表明该干预在此设置下有作用，但不能直接证明替换位置是唯一必要机制。替换值可能把系统推到自然数据从未访问的状态，多个冗余路径也会让单点消融效果较小。应比较多种扰动、对照位置和组合干预，并检查结论是否跨输入成立。

\section{稀疏自编码器与特征字典}

稀疏自编码器（sparse autoencoder, SAE）尝试把激活 $h\in\mathbb R^d$ 分解到更大的特征字典中。一个基本目标为
\[
z=\operatorname{ReLU}(W_eh+b_e),\quad \widehat h=W_dz+b_d,\qquad
\min\mathbb E\bigl[\|h-\widehat h\|_2^2+\lambda\|z\|_1\bigr].
\]
稀疏项鼓励每个输入只激活少量特征。为防止通过放大解码列、缩小编码系数任意降低 $L_1$ 惩罚，通常要约束解码字典列范数或使用等效尺度控制。字典学习研究展示了比单神经元分析更细的特征分解\cite{sae}。

重构误差、稀疏度、死特征比例和下游输出保持是不同指标。人为给某个特征命名只是解释假设，需要正反样本与干预证据支持。SAE 分解不一定唯一，也不能仅凭可解释特征数量增加就认定找到了模型的全部真实机制。

